\subsection{Reading a Dataset: From File to Query Pipeline}

When a Host wants to read an AnyBlox dataset, it first opens the \texttt{.any} file and reads its metadata during query planning. The metadata provides the Host with the schema, the total number of rows, and a recommended batch size, all of which are needed to partition the workload across threads and schedule the scan before decoding begins. The Host then passes this metadata and the \texttt{.any} file to \texttt{libanyblox}, which extracts the Wasm Decoder module, JIT-compiles it, and caches the result (subsequent uses of the same Decoder therefore incur no compilation cost). AnyBlox then sets up the execution environment for the Decoder and returns a job handle to the Host.

From this point, the Host repeatedly requests arbitrary tuple ranges from that handle. Modern query engines process data in batches to amortize operator communication overhead, so rather than decoding the entire dataset at once, the Host issues successive calls for ranges of tuples~\cite{kerstenEverythingYouAlways2018}. AnyBlox delegates each such request into the sandboxed Wasm module and returns the decoded result to the Host. The sole function through which this happens is \texttt{decode\_batch}~\cite{gienieczkoAnyBloxFrameworkSelfDecoding2025}:

\begin{itemize}
    \item \texttt{i32 data}: pointer to the start of the encoded data in Wasm linear memory
    \item \texttt{i32 data\_length}: length of the encoded data in bytes
    \item \texttt{i32 start\_tuple}: ID of the first tuple to decode
    \item \texttt{i32 tuple\_count}: number of tuples to decode
    \item \texttt{i32 state}: pointer to the state page in Wasm linear memory
    \item \texttt{i64 projection\_mask}: bitmask specifying which columns to decode, allowing the Host to skip columns that are not needed for the current query
\end{itemize}

The \texttt{data} parameter points to the encoded dataset inside the Decoder's execution environment, but how does the dataset get there in the first place? A naive solution would be to copy the entire dataset into a memory region accessible to the Decoder, but this would incur a heavy initialization cost proportional to the dataset size. Instead, AnyBlox maps the dataset file directly into the Decoder's linear memory via \texttt{mmap}, a system call that maps a file directly into the virtual address space of a process without copying it, so physical memory is only allocated when a page is first accessed.

Wasm's linear memory is a contiguous, byte-addressable memory region starting at address \texttt{0x0} that represents the entire address space accessible to the Decoder, isolated from the rest of the host process. AnyBlox controls its layout: the encoded dataset is mapped into it via \texttt{mmap}, and the state page is placed immediately after. The \texttt{data} parameter is therefore simply a pointer to the start of the mapped dataset within this region. AnyBlox reuses the linear memory region across successive jobs on the same thread, avoiding repeated initialization costs.

The \texttt{state} parameter points to a page within the linear memory that is zero-initialized at the start of every job. A Decoder can use it to cache metadata or memory allocations across successive \texttt{decode\_batch} calls. For example, a run-end encoding Decoder that performs a binary search on the first call can store its current position on the state page and resume directly on the next call, skipping the search entirely. Because the state page resets at the start of each job, Decoders are idempotent: identical parameters on the same dataset always produce equivalent results.

Once decoding is complete, \texttt{decode\_batch} returns a pointer to an Apache Arrow Record Batch within the linear memory. Arrow\footnote{Apache Arrow, \url{https://arrow.apache.org/}, (accessed: June 2, 2026)} is a columnar in-memory format developed by the Apache Software Foundation: rather than storing data row by row, it organizes each column as a contiguous array in memory, which allows query engines to process large amounts of data efficiently without unnecessary memory accesses. A Record Batch is a collection of such arrays, all of the same length, representing a horizontal slice of a table. Arrow standardizes the representation of common SQL types (e.g., integers, floating-point numbers, strings, dates, and more) and is supported by a wide range of query engines. The choice of Arrow as the uniform output format of AnyBlox follows the same principle as other intermediate representations in systems design, because by agreeing on a common language between two components, the number of required translations reduces from $N \times M$ to $N + M$. Every Decoder, regardless of the underlying encoding scheme, produces an Arrow Record Batch, and every Host that speaks Arrow can consume it without any AnyBlox-specific logic. AnyBlox translates the Wasm-internal pointer into a native address and passes the batch to the Host, where it enters the query pipeline. The degree of integration varies by system: DataFusion~\cite{lambApacheArrowDataFusion2024} uses Arrow as its primary internal representation throughout, making the transfer zero-copy end-to-end. DuckDB~\cite{raasveldtDuckDBEmbeddableAnalytical2019} and Spark\footnote{Apache Spark, \url{https://spark.apache.org/} (accessed: June 2, 2026)} do not use Arrow internally, but both have built-in support for reading Arrow batches and converting them into their respective internal representations on the fly. This conversion logic is part of the mainline code of both systems and requires no AnyBlox-specific adaptation~\cite{gienieczkoAnyBloxFrameworkSelfDecoding2025}.